\section{Deep Learning Fundamentals}
\label{section:primerdl}

\subsection{Artificial Neural Networks}

\begin{TCBBox}[label=Box:supervisedlearning]{Basic Introduction to Supervised Learning}

Deep learning algorithms have most commonly been adapted for causal inference using supervised machine learning, the most popular learning framework within the field.\footnote{The other two prominent paradigms are unsupervised learning and reinforcement learning.} The goal of supervised learning is teach a model a non-linear function that transforms covariates/features $X$ into predicted outcomes $\hat{Y}$ in unseen data. The model learns this function from labeled examples of $X_\text{tr}$ and $Y_{tr}$ in a \textbf{training dataset}.

As in traditional statistical analyses, the function is learned by optimizing the model's parameters such that they minimize the error between its predictions $\hat{Y_{tr}}$ and the true values $Y_{tr}$ using a \textbf{loss function} (e.g., a likelihood). In a traditional social science analysis focused on inference, we would stop here and interpret these parameters. In supervised machine learning where the focus is on generalization to unseen data, the model is ultimately used to predict outcomes $Y_{te}$ in a \textbf{test dataset} of previously unseen covariates/features $X_{te}$. This framework can be generically applied to cases where $Y$ is categorical (called classification problems), and where $Y$ is continuous (called regression problems).

Statistical learning theory articulates the central challenge of supervised learning as a balance between \textbf{overfitting} and \textbf{underfitting} the training dataset. This is also called the \textbf{bias-variance" tradeoff}. In a regression context, bias error is the difference between the expected value of Y and the expected value of the mapping function learned by the model.\footnote{Note that bias in statistical learning theory is not equivalent to bias of a statistical estimator.} High bias typically results from an algorithm that has not sufficiently learned the relationships in the training dataset (i.e., underfit the data). In contrast, an algorithm that has learned the training dataset so closely that it is fitting noise in the sample (i.e., overfitting) is likely to generalize poorly, producing out-of-sample predictions with high variance. Underfitting can be easily diagnosed and addressed by increasing the complexity of the model. In the case of deep learning, model complexity can be increased by adding additional layers or parameters/neurons.

Diagnosing and addressing overfitting is a more challenging problem. In supervised learning, overfitting is diagnosed after training (but before testing) by assessing predictive performance in a reserved portion of the training set called the \textbf{validation set}. If the model fits the training dataset well but performs poorly in the validation set, it is likely to generalize poorly to the test set as well. To prevent overfitting, \textbf{regularization} techniques can be used to simplify the complexity of the model. Training and regularization of neural networks is discussed in detail in Section \ref{section:practice}. For a full treatment of supervised learning and statistical learning theory, see \citet{hastie2009elements}.
\end{TCBBox}

Artificial neural networks (ANN) are statistical models inspired by the human brain \citep{brand_koch_xu,Goodfellow-et-al-2016}. In an ANN, each ``neuron" in the network takes the weighted sum of its inputs (the outputs of other neurons) and transforms them using a differentiable, non-linear function (e.g. sigmoid, rectified linear unit) that outputs a value between 0 and 1 if the transformed value is above some threshold. Neurons are arrayed in layers where an input layer takes the raw data, and neurons in subsequent layers take the weighted sum of outputs in previous layers as input. An ``output" layer contains a neuron for each of the predicted outcomes with transformation functions appropriate to those outcomes. For example, a regression network that predicts one outcome will have a single output neuron without a transformation function so that it produces a real number. A regression network without any hidden layers corresponds exactly to a generalized linear model (Fig. \ref{fig:perceptron}A). When additional ``hidden" layers are added between the input and output layers, the architecture is called a \textbf{feed-forward network} or \textbf{multi-layer perceptron} (Fig. \ref{fig:perceptron}B). A neural network with multiple hidden layers is called a ``deep" network, hence the name ``deep learning" \citep{lecun2015deep}. A neural network with a single, large enough hidden layer can theoretically approximate any continuous function \citep{Cybenko1989ApproximationBS}. 

\begin{TCBBox}[label=Box:training]{Reading Machine Learning Papers: Computational Graphs and Loss Functions}
Within the machine learning literature, novel algorithms are often presented in terms of their computational graph and loss function. A computational graph (not to be confused with a causal graph) uses arrows to depicts the flow of data from the inputs of a neural network,  through parameters, to the outputs. Layers of neurons or specialized sub-architectures are often generically abstracted as shapes. In our diagrams, we use rounded purple shapes to represent observables, orange rectangles for representation layers of the network, rounded white shapes for produced outputs, and textured rectangles for outcome modeling layers. Operations that are computed \textit{after} prediction (i.e., for which an error gradient is not calculated) are shown with dashed lines (e.g., plug-in estimation of causal estimands).

Along with the architecture, the loss function of a neural network is the primary means for the analyst to dictate what types of representations a neural network learns and what types of outputs it produces. In multi-task learning settings, we denote joint loss functions for an entire network as a weighted sum of the losses for substituent tasks and modules. These specific losses are weighted by hyperparameters. For example, we might weight the joint loss for a network that predicts outcomes and propensity scores as:

\begin{equation*}
    \argmin_{h,\pi}\mathcal{L} = \mathcal{L}_h + \lambda \mathcal{L}_\pi = MSE(Y,h(X,T)) + \lambda BCE(T,\pi(X,T))
\end{equation*}
where $h(X,T)$ is the predicted potential outcome, $\pi(X,T)$ is the predicted propensity score, $\lambda$ is a hyperparameter and MSE and BCE stand for mean squared error and binary cross entropy (i.e., log loss), common losses for regression and binary classification respectively (Box \ref{Box:Notation}). 
%\end{TCBBox}
\end{TCBBox}

\begin{figure}
    \centering
    \includegraphics[width=\linewidth]{content/Perceptronv2.png}
    \caption{\textbf{A: Generalized linear model represented as a computational graph.} Observable covariates $X_1,X_2,X_3$ and treatment status $T$ depicted in rounded purple boxes. Each of the lines between the  rounded purple inputs and the textured box represents a parameter (i.e., a $\beta$ in a generalized linear model equation). The textured box is an ``output neuron" that sums it's weighted inputs, performs a transformation $g$ (the link function in GLM; in this case the identity function), and predicts the conditional outcome $\hat{Y}(T)$. Instead of theoretically interpreting these parameters from an inferential statistics perspective, machine learning approaches typically use the predicted observed and unobserved potential outcomes for plug-in estimation of causal estimands (e.g., the conditional average treatment effect $\hat{CATE}$). After training, setting $T$ to $1-T$ for each observation can predict the unobserved potential outcome $\hat{Y}(1-T)$. Because this operation occurs after prediction and does not feed a gradient back to the network to optimize the parameters, it is depicted here with a dotted line. Plug-in calculation of $\hat{CATE}$ is similarly shown with a dotted line.
    \\
    \textbf{B: Feed-forward neural network (S-learner).} In a feed-forward neural network, additional fully connected (parameterized) layers of neurons are added between the inputs (rounded purple) and output neuron. The size of the input covariates and hidden layers are generically abstracted as boxes (orange). The final hidden layer before the output neuron is denoted $\Phi$ because the hidden layers collectively encode a representation function (see section \ref{section:replearning}). In causal inference settings, this architecture is sometimes called a S(ingle)-learner because one feed-forward network learns to predict both potential outcomes.}
    \label{fig:perceptron}
\end{figure}

Neural networks are trained to predict their outcomes by optimizing a \textbf{loss function} (also called an objective or cost function). During training, the \textbf{backpropagation} algorithm
uses calculus's chain rule to assign portions of the total error in the loss function to each neuron in the network. An optimizer, such as the stochastic gradient descent algorithm or the currently popular ADAM algorithm \citep{Kingma2014AdamAM}, then moves each parameter in the opposite direction of this error gradient. Neural networks first rose to popularity in the 1980s but fell out of favor compared to other machine learning model families (e.g., support vector machines) due to their expense of training. By the late 2000s, improvements to backpropagation, advances in computing power (i.e., graphic cards), and access to larger datasets collectively enabled a deep learning revolution where ANNs began to significantly outperform other model families. Today, deep learning is the hegemonic machine learning approach in industries and fields other than social science.

\subsection{Deep Learning in Practice}


\label{section:practice}
This section focuses on the practice of training neural networks within a supervised learning framework. While the principles behind supervised machine learning are universal, the workflow for neural networks differs substantially from other ML approaches (e.g., random forests, support vector machines) in practice. Figure \ref{fig:workflow} presents this workflow in four different parts: Set Up, Training, Model Evaluation, and Interpretation. We delve into each of these topics in more detail below. Box \ref{Box:supervisedlearning} contains a basic introduction to supervised learning for unfamiliar readers.

\begin{figure}
    \centering
    \includegraphics[width=\linewidth]{content/ModelSplittingCriteria.png}
    \caption{\textbf{Supervised Deep Learning Workflow.} \textbf{1) Set Up:} The first step in training a deep learning model is splitting the data into a training set, validation set, and optionally a test set. Initial hyperparameters are then selected from a set of choices specified by the user. \textbf{2) Training:} In each iteration of the training process (called an epoch), the training set is randomly divided into small minibatches For each minibatch, the network makes predictions for all units, and calculates the error gradients to be assigned to each neuron in the network based on those predictions. An optimizer then move the network's parameters in the opposite direction of the error gradient. After all minibatches have been trained (one epoch), error is calculated on the entire validation set. This whole process is repeated up until the validation error stops decreasing (to avoid overfitting). \textbf{3) Model Evaluation:} A criterion (typically the validation error) is used to evaluate the performance of this hyperparameterization. New hyperparameters are then selected using a hyperparameter optimization algorithm (eg. Grid search, Bayesian hyperparameter optimization, genetic algorithms) and steps 1 and 2 are repeated. Once the hyperparameter optimization algorithm has completed its search, the ``best" model is selected for inference. \textbf{4) Inference and interpretation:} With a model selected, the analyst is now ready to apply it to their test data (or in the case of statistical inference, potentially the full dataset). Predictions of the outcomes and/or propensity score can then be used to compute the CATE and calculate confidence intervals. Feature importance algorithms like SHAP or Integrated Gradients can also be used to interpret the CATE estimates.}
    \label{fig:workflow}
\end{figure}


\subsubsection{Set Up and Hyperparameters}

The first step in training a neural network, as in other types of supervised machine learning, is to split your dataset into training, validation, and testing datasets (Fig. \ref{fig:workflow}A). If the network is being used for statistical inference, as here, the testing dataset is optional, and inference may be conducted on just the validation set or the full dataset.

While the computational graph and loss function define a deep learning architecture (Box \ref{Box:training}), actual implementations can vary significantly due to the choice of hyperparameters. In supervised machine learning, \textbf{hyperparameters} are parameters that are not learned automatically when training the model, but must be specified by the analyst. In deep learning, architectural hyperparameters include the number of layers to use for each section of the computational graph, the number of neurons to use in each layer, and the activation functions to be used by neurons. While some basic rules of thumb apply (e.g., use fewer layers than neurons), these choices remain poorly understood theoretically\footnote{For some interesting work on understanding neural networks theoretically from a statistical physics perspective see \citet{PDLT-2022}.}; Decisions are generally made by comparing empirical performance on the validation set, a practice called \textit{hyperparameter tuning}.

\subsubsection{Training and Regularization}
\label{subsection:trainingregularization}

Neural networks are trained by repeatedly making predictions from the training set, calculating error gradients for each parameter, and backpropagating small fractions of those error gradients. (Fig. \ref{fig:workflow} B).  A full pass through examples in the training set is called a training loop or \textbf{epoch}. At the beginning of each epoch, the training set is divided into \textbf{mini-batches} of 2 to 1000 units, randomly sampled without replacement. This practice not only aids in memory management, it also improves optimization. Using small random samples reduces the risk of large “exploding” error gradients, particularly early in the training, that could cause the model to overshoot optimal solutions and instead get stuck in local minima.

The size of mini-batches can be considered a hyperparameter.\footnote{In the specific context of causal inference, we recommend not having mini-batches that are too small such that the model can learn from both treated and control units with sufficient overlap.} Because a mini-batch of data is only a sample of a sample (the training dataset), the optimizer only adjusts weight parameters by a fraction of the error gradient (the \textbf{learning rate}) to avoid overfitting. The learning rate is also a hyperparameter, that typically varies between 0.0001 and 0.01.

The non-convex nature of most loss functions\footnote{In convex functions (e.g. the OLS loss), there is a single minimum, so optimizing the function means that you will always converge at the same parameter weights. This is not the case for non-convex functions which may have many local minima.} means that optimization often requires hundreds to potentially millions of epochs of training. Moreover, neural networks are highly susceptible to overfitting because it is easy to overparameterize them with excessive neurons/layers. To ward against overfitting, error metrics on the complete validation set are computed at the end of every epoch. In a regularization practice called “\textbf{early stopping},” analysts usually stop training once validation metrics stop improving.
Other common regularization techniques include \textbf{weight decay} (i.e., $\ell^2$ norm, ridge, or Tikhonov) penalties on the parameters, dropout of neurons during training, and batch normalization. 

\textbf{Dropout} is a regularization technique in deep learning where certain nodes are randomly silenced from training during a given epoch \citep{Srivastava2014DropoutAS}. The general idea of dropout is to force two neurons in the same layer to learn different aspects of the covariate/feature space and reduce overfitting. \textbf{Batch normalization} is another regularization technique applied to a layer of neurons \citep{Ioffe2015BatchNA}. By standardizing (i.e. z-scoring) the inputs to a layer on a per-batch basis and then rescaling them using trainable parameters, batch normalization smooths the optimization of the loss function. The addition and extent of each of these regularization techniques can be treated as hyperparemeters.

\subsubsection{Model Selection}
(\textit{Tutorial 2 \href{https://colab.research.google.com/drive/1MBUkQO7rh89JrlAV0p1aLNoF9CIB3rZZ?usp=sharing}{\includegraphics[height=\fontcharht\font`\B,keepaspectratio]{content/colab_button.png} }})
\label{subsection:modelselection}

After the model has been trained, the analyst compares models assembled with different hyperparameterizations or initial parameter values (Fig. \ref{fig:workflow}C). Hyperparameterizations can be chosen using random search, an exhaustive grid search of all possible combinations, or strategic search algorithms like Bayesian hyperparameter optimization or evolutionary optimization \citep{snoek2012}. Validation loss metrics on the final epoch are commonly used for these comparisons. 

Model selection for causal estimators is complicated by the fundamental problem of causal inference: we are not actually interested in the observed “factual” outcomes and propensity scores, but the CATE and ATE. In the case of algorithms like Dragonnet \ref{section:ipw} where the validation loss explicitly targets a causal quantity, we use that as the model selection criterion. In cases where the algorithm is only trained for outcome modeling or propensity modeling, other solutions are needed. In the Appendix, we describe \citet{johannson2020}'s proposal to use matching on a nearest neighbor approximation of the \textit{Precision in Estimated Heterogeneous Effects} (PEHE), a measure of CATE bias, as an alternative model selection metric (Appendix A.\ref{appendix:peheselection}).

The development of more sophisticated methods for model selection of causal estimators through data simulation is an active area of research within this literature.\footnote{We note that crossfitting \citep{zivich2021machine}, another approach that has emerged for model selection of other types of machine learning causal estimators may work for the models discussed here, but is likely data-inefficient.} For example, \citet{parikh2022} use deep generative models to approximate the data generating distribution under weak, non-parametric assumptions. \citet{alaa2019validating} independently model each  outcome and the propensity score before using influence functions to assess model error.

\subsection{Representation Learning and Multitask Learning}

\label{section:replearning}
One comparative advantage of deep learning over other machine learning approaches has been the ability of ANNs to encode and automatically compress informative features from complex data into flexible, relevant ``\textbf{representations}" or ``embeddings" that make downstream supervised learning tasks easier \citep{Goodfellow-et-al-2016, bengio2013deep}. While other machine learning approaches may also encode representations, they often require extensive pre-processing to create useful features for the algorithm (i.e., feature engineering). Through the lens of representation learning, a geometric interpretation of the role of each layer in a supervised neural network is to transform its inputs (either raw data or output of previous layers) into a typically lower (but possibly higher) dimensional vector space. As a means to share statistical power, encoded representations can also be jointly learned for two tasks at once in \textbf{multi-task learning}. 

The simplest example of a representation might be the final layer in a feed-forward network, where the early layers of the network can be understood as non-linearly encoding the inputs into an array of latent linear features for the output neuron \citep{Goodfellow-et-al-2016} (Fig. \ref{fig:perceptron}B). A famous example of representation learning is the use of neural networks for face detection. Examining the representations produced by each layer of these networks shows that each subsequent layer seems to capture increasingly abstract features of a face (first edges, then noses and eyes, and finally whole faces) \citep{lecun2015deep}. A more familiar example of representation learning to social scientists might be word vector models like Word2Vec \citep{mikolov2013}. Word2Vec is a neural network with one hidden layer and one output layer where words that are semantically similar are closer together in the representation space created by the hidden layer of the network. 

\emph{The novel contribution of deep learning to causal estimation is the proposal that a neural network can learn a function $\Phi$ that produces representations of the covariates decorrelated from the treatment.} Fundamentally, the idea is that $\Phi$ can transform the treated and control covariate distributions into a representation space such that they are indistinguishable (Fig. \ref{fig:balancing}). To ensure that these representations are also still predictive of the outcome (multi-task learning), multiple loss functions are generally applied simultaneously to balance these objectives. This approach is applied in a majority of the algorithms presented in section \ref{section:mainbody}.

\begin{figure}
    \centering
    \includegraphics[width=\linewidth]{content/balancing.png}
    \caption{\textbf{Balancing through representation learning.} The promise of deep learning for causal inference is that a neural network encoding function $\Phi$ can transform the treated and control covariate distributions into a representation space such that they are indistinguishable. Used with permission from \citet{shenjohannson2018}.}
    \label{fig:balancing}
\end{figure}

